\textbf{Problème 1.10 (Dujella, 2026).} Définir un analogue d'un $D(n)$-$m$-uplet dans un anneau non commutatif et démontrer des résultats non triviaux concernant l'existence de $D(n)$-quadruplets, par exemple dans l'anneau des matrices $2\times 2$ à coefficients entiers.

\medskip

\textit{Remarque.} Dujella et Franušić (arXiv:2604.23404) proposent de considérer le produit de Jordan $A \circ B = \frac{1}{2}(AB+BA)$. La notion de triplet régulier se généralise alors naturellement : si $A \circ B = R^2$, alors $C = A+B+2R$ vérifie $A \circ C = (A+R)^2$ et $B \circ C = (B+R)^2$.
